\chapter{Glossary}
\label{app:glossary}

\begin{description}[style=nextline,leftmargin=1.2em,itemsep=4pt]
\item[Absolute priority rule] The rule in cramdowns under Chapter 11 of the US Bankruptcy Code that a dissenting class must be paid in full before junior classes receive anything.
\item[Activity level proposition] A proposition about how a trace came to be where it was found, as distinct from whose trace it is.
\item[Adversarial procedure] A trial procedure in which the parties gather and present evidence and the judge acts as umpire.
\item[Balance of probabilities] The civil standard of proof: a fact is proved if it is more probable than not.
\item[Base rate] The frequency of an event or characteristic in a reference population, used to form a prior.
\item[Bayes factor] The ratio of the probabilities of the evidence under two hypotheses, averaged over any unknown parameters; the factor by which evidence multiplies prior odds.
\item[Bayes risk] The expected loss of a decision rule, averaged over states and data.
\item[Bayes' rule] The theorem that posterior probability is proportional to likelihood times prior.
\item[Bayesian network] A directed acyclic graph with conditional probability tables, representing a joint distribution and its conditional independences.
\item[Bellwether trial] A trial of a selected case in mass litigation, used to inform the valuation of others.
\item[Beyond reasonable doubt] The criminal standard of proof; in England, the jury must be sure of guilt.
\item[Burden of proof] The allocation to a party of the risk of losing if the evidence does not reach the standard.
\item[But for test] The test of causation that asks whether the harm would have occurred without the defendant's breach.
\item[Calibration] The property of a probability forecaster that events forecast with probability $p$ occur with frequency close to $p$.
\item[Chain event graph] A graphical model built from an event tree, suited to asymmetric scenarios.
\item[Collider] A node on a path at which two arrows meet, as in $A\to B\leftarrow C$; observing it, or one of its descendants, can make $A$ and $C$ dependent.
\item[Conditional independence] Two quantities are conditionally independent given a third if, once the third is known, learning one says nothing about the other.
\item[Conformal prediction] A method that produces prediction sets with guaranteed marginal coverage under exchangeability.
\item[Corroboration] Independent evidence supporting a proposition already supported by other evidence.
\item[Cost benefit analysis] The comparison of the expected costs and benefits of a regulation or decision, in money terms.
\item[Counterfactual] A statement about what would have happened under conditions different from those that held.
\item[Decibans] Units of weight of evidence: ten times the base ten logarithm of the likelihood ratio.
\item[Defence fallacy] The error of treating evidence as worthless because it does not single out the defendant.
\item[Deference] The weight a court gives to an agency's interpretation of a statute; after \emph{Loper Bright}, a matter of persuasive weight rather than a rule of decision.
\item[Dershowitz fallacy] The error of quoting a probability conditioned on less than what is known, such as quoting the probability that a batterer kills his partner when the partner is known to have been killed.
\item[Disparate impact] Liability for a neutral practice that disproportionately disadvantages a protected group without adequate justification.
\item[Event study] A statistical method for estimating the effect of news on a security's price.
\item[Expected utility] The probability weighted average of the utilities of possible consequences of an action.
\item[Explaining away] The reduction in the probability of one cause when another cause of the same effect is established.
\item[Fixed recoverable costs] A schedule of legal costs recoverable by a winning party, fixed in advance by rules of court.
\item[Fraud on the market] The presumption that investors rely on public misstatements through the market price.
\item[Generative Bayesian computation] Approximation of posterior distributions by training neural networks on simulated parameter and data pairs.
\item[Hand formula] The rule that a defendant who fails to take a precaution is negligent if its burden is less than the probability of harm times its gravity.
\item[Hierarchical model] A model in which parameters for groups are drawn from a common distribution, producing partial pooling.
\item[Hindsight bias] The tendency to judge an outcome as more foreseeable after it has occurred.
\item[Inquisitorial procedure] A procedure in which a judge directs the investigation and gathers evidence.
\item[Intime conviction] The inner conviction of the judge or jury, the standard of French criminal procedure.
\item[Island problem] The model of forensic identification in which a criminal is one of a finite population, some of whom share a characteristic.
\item[Kelly criterion] The rule for sizing a sequence of independent bets that maximises the long run growth rate of wealth.
\item[Likelihood ratio] The probability of evidence under one hypothesis divided by its probability under another.
\item[Lindley's paradox] The phenomenon that a result just significant at a fixed level gives increasing Bayesian support to a point null hypothesis, against a diffuse alternative, as sample size grows.
\item[Litigation funding] Payment of a party's legal costs by a third party in return for a share of the recovery.
\item[Loss function] A function giving the loss from each action in each state of the world.
\item[Loss of a chance] A head of damages for the reduction in the probability of a favourable outcome.
\item[Marginal likelihood] The probability of the data under a model, averaged over its parameters.
\item[Material contribution] A basis of causation where the defendant's breach materially contributed to the harm or the risk of it.
\item[Multiplier] In English damages for future loss, the factor applied to an annual loss to give a lump sum, reflecting duration and discounting.
\item[Naked statistical evidence] Evidence consisting only of a base rate, without individual evidence about the case.
\item[NESS test] The test that treats a condition as a cause if it was a necessary element of a sufficient set.
\item[Noise] Unwanted variability in judgments that should be the same.
\item[Objective] The aim a decision serves, expressed through a utility or loss function.
\item[Odds] The ratio of the probability of a proposition to the probability of its negation.
\item[Open texture] The indeterminacy of general terms in penumbral cases.
\item[Periodical payments] Damages paid in instalments over the claimant's life rather than as a lump sum.
\item[Plausibility pleading] The American standard requiring a complaint to state a claim that is plausible on its face.
\item[Posterior] The probability of a proposition after taking account of evidence.
\item[Power] The probability that a test or procedure correctly rejects a false null hypothesis; in criminal terms, the probability of convicting the guilty.
\item[Price impact] The effect of a misstatement or disclosure on a security's price.
\item[Prior] The probability of a proposition before taking account of the evidence in question.
\item[Probability of causation] The probability that an exposure caused a harm in an individual who was exposed and harmed.
\item[Proper scoring rule] A rule for rewarding probability forecasts that is maximised in expectation by honest reporting.
\item[Proportionate remedy] A remedy scaled to a probability or to the ratio of the actual to the counterfactual position.
\item[Prosecutor's fallacy] The error of reading the probability of the evidence given innocence as the probability of innocence given the evidence.
\item[Random match probability] The probability that a person unrelated to the source shares a forensic profile.
\item[Regression discontinuity] A design that estimates effects by comparing units just on either side of a threshold or border.
\item[Relevance] The property of evidence whose likelihood ratio differs from one.
\item[ROC curve] A curve showing the trade off between the Type I error rate and power as a decision threshold varies.
\item[Selection effect] Bias arising because the cases observed are not a random sample of those of interest.
\item[Standard of proof] The posterior probability a fact finder must reach before acting on a fact.
\item[Stare decisis] The doctrine that courts follow precedent.
\item[Type I error] Rejecting a true null hypothesis; in criminal law, convicting an innocent person.
\item[Type II error] Failing to reject a false null hypothesis; in criminal law, acquitting a guilty person.
\item[Value of information] The expected improvement in a decision from observing further evidence before deciding.
\item[Weight of evidence] The logarithm of the likelihood ratio, in this book ten times its base ten logarithm, measured in decibans; additive for conditionally independent items.
\item[Wigmore chart] A diagram of the inferential relations among items of evidence and propositions.
\end{description}
